\begin{definition}
   Se llama \textit{función} a la regla que asigna a cada elemento de un conjunto
    $A$  un y s\'olo un elemento de un conjunto $B$. Tomaremos, en esta secci\'on,  el caso  donde $A, B \subseteq \mathbb{R}.$   La manera habitual de denotar una funci\'on $f$ es $$f: A \subseteq \mathbb{R} \rightarrow B \subseteq \mathbb{R}.$$  A los conjuntos $A$ y $B$  se los llama, respectivamente, \textit{dominio}  y \textit{codominio} de la funci\'on y  la notaci\'on usual para los mismos  es  $A=\text{Dom}(f)$ y $B=\text{Cod}(f)$.
    \label{def:FuncMate1}
\end{definition}


\begin{definition} Dada una funci\'on $f: A \subseteq \mathbb{R} \rightarrow B \subseteq \mathbb{R},$  se llama \textit{imagen} de una función y se nota  $\text{Im}(f)$, a
    al conjunto 
    \begin{equation*}
        \text{Im}(f) = \left\{ y\in B: \exists\;x\in A\:|\:f(x)=y \right\}.    
    \end{equation*}
\end{definition}

\begin{definition}\label{grafm1}
   Dada una función $f: A \subseteq \mathbb{R} \rightarrow B \subseteq \mathbb{R}$ se define la \textit{gráfica} de $f$, y se nota  $\text{Graf}(f)$, al conjunto
    \begin{equation*}
        \text{Graf}(f) = \left\{ (x,y)\in \mathbb{R}^{2}: x \in A \land y=f(x)\right\}.    
    \end{equation*}
\end{definition}


%\textcolor{red}{Aca hay un lio en los ejemplos, la tabla y el grafico no corresponden a la formula. Hacer dos ejemplos completo por separado. }

\begin{example}
    Sea $$f:  [0,\infty]  \subseteq \mathbb{R} \rightarrow  \mathbb{R} \:|\:  f(x)=2x-1$$ tenemos que, $f$ es el nombre de la función, 
    $x$ es la variable independiente de $f$  y $2x-1$ su  fórmula, es decir,  el valor que la asigna $f$ a cada $x \in  [0,+\infty). $ 
    Adem\'as,   $\text{Dom}(f)=  [0,+\infty)$,   $\text{Cod}(f)=  \mathbb{R}$ y $\text{Im}(f)= [-1,+\infty)$ .
    Si se quisiera, se puede hacer una tabla de valores para la función:
    \begin{table}[H]
        \centering
        \begin{tabular}{|c|c|}
            \hline
            $x$ & $y=f(x)=2x-1$ \\ \hline
            -2  & -5        \\ \hline
            -1  & -3        \\ \hline
            0   & -1        \\ \hline
            1   & 1        \\ \hline
            2   & 3        \\ \hline
            \end{tabular}
      \end{table}

      Estos puntos, se pueden graficar sobre un plano cartesiano.
      \begin{figure}[H] 
        \begin{center}
            \begin{tikzpicture}
                \begin{axis}[
                    axis lines = middle,
                    xmin = -2.5, xmax = 2.5,
                    ymin = -5.5, ymax = 3.5,
                    xlabel = {$x$},
                    ylabel = {$y$},
                    domain = -2:2,
                    samples = 100,
                    xtick={-2,-1,1,2},
                    ytick={1,2,3,4},
                    clip=false,
                    ]
                
                    % Define the curves
                    \fill[red] (0,-1) circle (4pt);
                    \fill[red] (-2,-5) circle (4pt);
                    \fill[red] (-1,-3) circle (4pt);
                    \fill[red] (1,1) circle (4pt);
                    \fill[red] (2,3) circle (4pt);
                \end{axis}
            \end{tikzpicture}
        \end{center}
        \caption{Conjunto de puntos pertenecientes a la tabla de valores de $f(x)=2x-1$ posicionados sobre el plano cartesiano.}
        %\label{fig:unitCircle1} % Etiqueta para hacer referencia a la imagen
    \end{figure}
    Además, sabemos que la función corresponde a una función lineal, por lo que, para graficar los infinitos puntos pertenecientes
    a la $\text{Graf}(f)$ podemos unir los anteriores con una recta.
    \begin{figure}[H] 
        \begin{center}
            \begin{tikzpicture}
                \begin{axis}[
                    axis lines = middle,
                    xmin = -2.5, xmax = 2.5,
                    ymin = -5.5, ymax = 3.5,
                    xlabel = {$x$},
                    ylabel = {$y$},
                    domain = -2:2,
                    samples = 100,
                    xtick={-2,-1,1,2},
                    ytick={1,2,3,4},
                    clip=false,
                    ]
                
                    % Define the curves
                    \addplot[name path=A, blue, thick, domain=-2:2] {2*x-1};
                    \fill[red] (0,-1) circle (4pt);
                    \fill[red] (-2,-5) circle (4pt);
                    \fill[red] (-1,-3) circle (4pt);
                    \fill[red] (1,1) circle (4pt);
                    \fill[red] (2,3) circle (4pt);
                \end{axis}
            \end{tikzpicture}
            \end{center}
        \caption{En rojo, los puntos de la tabla de valores. En azul, $\text{Graf}(f)$, con sus infinitos puntos, en el intervalo $[-2,2]$.}
        %\label{fig:unitCircle1} % Etiqueta para hacer referencia a la imagen
    \end{figure}

  \end{example}
\begin{example}  Sea $$f: [-2,2]  \subseteq \mathbb{R} \rightarrow  \mathbb{R} \:|\:  f(x)=x^2$$ 

 Se arma la siguiente tabla de valores, para valores discretos de $x\in [0,2]$:
    \begin{table}[H]
        \centering
        \begin{tabular}{|c|c|}
            \hline
            $x$ & $y=f(x)=x^2$ \\ \hline
            -2  & 4        \\ \hline
            -1  & 1        \\ \hline
            0   & 0        \\ \hline
            1   & 1        \\ \hline
            2   & 4        \\ \hline
            \end{tabular}
      \end{table}
    En este caso, haciendo un análisis de la función, se puede ver que 
    $\text{Dom}(f)=  [-2,2]$,   $\text{Cod}(f)=  \mathbb{R}$ y $\text{Im}(f)= [0,4]$
    Luego, el siguiente conjunto de puntos, se encuentran incluidos en la
    gráfica de la función (que cabe aclarar, tiene infinitos elementos).
    \begin{equation*}
        \left\{ (-2,4),(-1,1),(0,0),(1,1),(2,4) \right\}\subset \text{Graf}(f)
    \end{equation*}
    Este conjunto se puede graficar en el plano coordenado x-y de la siguiente forma.
    \begin{figure}[H] 
        \begin{center}
            \begin{tikzpicture}
                \begin{axis}[
                    axis lines = middle,
                    xmin = -2.5, xmax = 2.5,
                    ymin = 0, ymax = 4.5,
                    xlabel = {$x$},
                    ylabel = {$y$},
                    domain = -2:2,
                    samples = 100,
                    xtick={-2,-1,1,2},
                    ytick={1,2,3,4},
                    clip=false,
                    ]
                
                    % Define the curves
                    \fill[red] (0,0) circle (4pt);
                    \fill[red] (-2,4) circle (4pt);
                    \fill[red] (-1,1) circle (4pt);
                    \fill[red] (1,1) circle (4pt);
                    \fill[red] (2,4) circle (4pt);
                \end{axis}
            \end{tikzpicture}
        \end{center}
        \caption{Subconjunto de los puntos pertenecientes a la $\text{Graf}(f)$ posicionados sobre el plano x-y.}
        %\label{fig:unitCircle1} % Etiqueta para hacer referencia a la imagen
    \end{figure}
    
  
  \textcolor{red}{hacer dos dibujos, el primero solo con los puntos rojos y el segundo con los mismos puntos mas el trazo. Decir algo asi como: por conocimiemto previo sabemos que la manera de unir los puntos es a trav\'es de..... }
  
    Por conocimiento previo, sabemos que la forma de unir estos puntos para graficar $f(x)$ es a través de una parábola. Así, 
    graficando la función como uno lo haría normalmente, en realidad
    se están dibujando los infinitos puntos pertenecientes
    a la $\text{Graf}(f)$ sobre el plano.
        \begin{figure}[H] 
            \begin{center}
                \begin{tikzpicture}
                    \begin{axis}[
                        axis lines = middle,
                        xmin = -2.5, xmax = 2.5,
                        ymin = 0, ymax = 4.5,
                        xlabel = {$x$},
                        ylabel = {$y$},
                        domain = -2:2,
                        samples = 100,
                        xtick={-2,-1,1,2},
                        ytick={1,2,3,4},
                        clip=false,
                        ]
                    
                        % Define the curves
                        \addplot[name path=A, blue, thick, domain=-2:2] {x^2};
                        \fill[red] (0,0) circle (4pt);
                        \fill[red] (-2,4) circle (4pt);
                        \fill[red] (-1,1) circle (4pt);
                        \fill[red] (1,1) circle (4pt);
                        \fill[red] (2,4) circle (4pt);
                    \end{axis}
                \end{tikzpicture}
                \end{center}
            \caption{En rojo, los puntos del subconjunto de la gráfica de la figura anterior. En azul, $\text{Graf}(f)$, con sus infinitos puntos, en el intervalo $[-2,2]$.}
            %\label{fig:unitCircle1} % Etiqueta para hacer referencia a la imagen
        \end{figure}
\end{example}
\begin{definition}
  
  Sea la funci\'on $f: A \subseteq \mathbb{R} \rightarrow B \subseteq \mathbb{R}$. Diremos que
   $f$ es \textit{sobreyectiva} si se cumple que su codominio es igual a su imagen, esto es
   \begin{equation*}
        \text{Cod}(f)=\text{Im}(f).
    \end{equation*}   Es decir que cada elemento del conjunto de salida es la evaluación de la función en al menos un elemento del dominio.

 Diremos que $f$ es   \textit{inyectiva} si a cada elemento en el conjunto imagen le corresponde un único
    valor del  conjunto dominio, esto es,   si  $a,b\in \text{Dom}(f)$, entonces:
     \begin{equation*}
        a=b \iff f(a)=f(b),
    \end{equation*} 

    Finalmente, diremos que $f$  es \textit{biyectiva} cuando es simultáneamente \textit{sobreyectiva} e \textit{inyectiva} a la vez.
\end{definition}

\begin{definition}
    Sea $f(x)$ definida en un intervalo abierto alrededor de $x_0$, excepto posiblemente $x_0$, es decir, en un entorno reducido aldededor del mismo
    $\varepsilon_0(x_0)$, se define el \textit{límite de $f(x)$ cuando x se aproxima a $x_0$}:
    \begin{equation*}
        \lim_{x\to x_0}{f(x)}=L 
    \end{equation*}
    si y solo si
    \begin{equation*}
        \forall \epsilon>0 : \exists \delta >0 : 0<\left\lvert x-x_0 \right\rvert<\delta \Rightarrow \left\lvert f(x) - L\right\rvert < \epsilon 
    \end{equation*}
    \label{def:Limite}
\end{definition}

\begin{obs}{Límites Notables}
    \label{obs:LimitesNotables}
    A continuación se presenta una lista de límites notables que pueden ser de utilidad:

    \begin{minipage}{0.45\textwidth}
        \begin{equation}
            \lim_{x\to 0}{\frac{\sin{x}}{x}}=1
            \label{limNot:sen(x)/x}
        \end{equation}
        \begin{equation}
            \lim_{x\to 0}{\frac{1-\cos{x}}{x}}=0
            \label{limNot:(1-cos(x))/x}
        \end{equation}
        \begin{equation}
            \lim_{x\to 0}{\frac{\exp{kx}-1}{x}}=k
            \label{limNot:(e^{kx}-1)/(x)}
        \end{equation}
    \end{minipage}
    \hfill
    \begin{minipage}{0.45\textwidth}
        \begin{equation}
            \lim_{x\to 0}{\frac{a^{kx}-1}{x}}=k\ln{a}
            \label{limNot:(a^(kx)-1)/(x)}
        \end{equation}
        \begin{equation}
            \lim_{x\to +\infty}{\left( 1+\frac{1}{x} \right)^x}=e
            \label{limNot:e}
        \end{equation}
        \begin{equation}
            \lim_{x\to 0}{\ln{\left(\frac{1+x}{x} \right)}}=1
            \label{limNot:ln(1+x/x)}
        \end{equation}
    \end{minipage}
\end{obs}

\begin{definition}{\textbf{Polinomio de Taylor}}
    Sea $f$ una función con derivades de orden $n$ en $x=a \in \text{Dom}(f)$. Se llama \textit{Polinomio de Taylor} 
    de grado $n$ de la función $f$ en $x=a$, y se simboliza $T_{f,a,n}(x)$, a:
    \begin{equation*}
        T_{f,a,n}(x) = \sum^n_{k=0} \frac{f^{(k)}(a)}{k!}(x-a)^k 
    \end{equation*}
    \label{def:PolTaylor}
\end{definition}